% =====================================================================
% Parámetros de calificación — Sistema de scoring de atenciones
% Compilar:  latexmk -pdf parametros-de-calificacion.tex
%
% Este documento describe el sistema TAL COMO ESTÁ IMPLEMENTADO en
% 2026.08-rubricas-v24. Cada umbral corresponde a una constante del
% código y se cita el módulo donde vive.
% =====================================================================
\documentclass[11pt,a4paper]{article}

\usepackage[utf8]{inputenc}
\usepackage[T1]{fontenc}
\usepackage[spanish,es-nodecimaldot]{babel}
\usepackage{lmodern}
\usepackage[margin=2.3cm,headheight=14pt]{geometry}
\usepackage{xcolor}
\usepackage{booktabs}
\usepackage{tabularx}
\usepackage{longtable}
\usepackage{array}
\usepackage{enumitem}
\usepackage{titlesec}
\usepackage{fancyhdr}
\usepackage{tcolorbox}
\usepackage[hidelinks]{hyperref}

\definecolor{primary}{HTML}{1F3864}
\definecolor{accent}{HTML}{2E75B6}
\definecolor{soft}{HTML}{EEF3FA}
\definecolor{warn}{HTML}{C0504D}
\definecolor{ok}{HTML}{2E7D32}
\definecolor{gray70}{HTML}{4D4D4D}

\hypersetup{
  pdftitle={Parametros de calificacion - Sistema de scoring},
  pdfsubject={Especificacion de rubricas, umbrales y precedencias},
  pdfkeywords={scoring, rubricas, ATC, calificacion, umbrales}
}

\pagestyle{fancy}
\fancyhf{}
\fancyhead[L]{\small\color{gray70}Parámetros de calificación}
\fancyhead[R]{\small\color{gray70}\thepage}
\renewcommand{\headrulewidth}{0.4pt}

\titleformat{\section}{\normalfont\Large\bfseries\color{primary}}{\thesection}{0.6em}{}
\titleformat{\subsection}{\normalfont\large\bfseries\color{accent}}{\thesubsection}{0.6em}{}
\setlist[itemize]{leftmargin=1.3em,itemsep=2pt,topsep=3pt}
\setlist[enumerate]{leftmargin=1.5em,itemsep=2pt,topsep=3pt}

\newcommand{\cod}[1]{\texttt{\small #1}}
\newtcolorbox{nota}[1][]{colback=soft,colframe=accent,boxrule=0.6pt,
  arc=2pt,left=6pt,right=6pt,top=5pt,bottom=5pt,#1}
\newtcolorbox{alerta}[1][]{colback=white,colframe=warn,boxrule=0.8pt,
  arc=2pt,left=6pt,right=6pt,top=5pt,bottom=5pt,#1}

\newcommand{\est}[1]{\textbf{#1}\,\raisebox{0.5pt}{$\star$}}

\begin{document}

\begin{titlepage}
\centering
\vspace*{3cm}
{\Huge\bfseries\color{primary} Parámetros de calificación\par}
\vspace{0.6cm}
{\Large\color{gray70} Sistema de scoring de atenciones\par}
\vspace{2cm}
{\large Versión del scoring: \cod{2026.08-rubricas-v24}\par}
\vspace{0.4cm}
{\large 31 de agosto de 2026\par}
\vfill
\begin{minipage}{0.82\textwidth}
\small\color{gray70}
Este documento describe el sistema \textbf{tal como está implementado}. Cada umbral
que aparece en él es una constante del código, y se indica el módulo donde vive, para
que cualquier afirmación se pueda verificar contra la fuente. No describe cómo
debería funcionar: describe cómo funciona hoy.
\end{minipage}
\vspace{2cm}
\end{titlepage}

\tableofcontents
\newpage

% =====================================================================
\section{Cómo se lee una nota}

\subsection{La escala}

La escala es \textbf{una sola para todos los motivos}, y la definió el negocio el
6 de agosto de 2026.

\begin{center}
\begin{tabular}{@{}clL@{}}
\toprule
\textbf{Estrellas} & \textbf{Etiqueta} & \textbf{Qué significa} \\
\midrule
5 & excelente  & Se logró el \textsc{mejor escenario} del motivo \\
4 & buena      & Se hizo bien \\
3 & aceptable  & Faltó algo leve \\
2 & deficiente & Faltaron varias cosas \\
1 & mala       & Se demoró mucho \emph{y} contestó mal, o no contestó \\
\bottomrule
\end{tabular}
\end{center}

\begin{nota}
\textbf{La estrella no la elige el modelo.} El LLM juzga \emph{hechos} (¿atendió el
motivo?, ¿hizo algo extra?) y el \emph{código} traduce esos hechos a una etiqueta, y
la etiqueta a una estrella, con una tabla fija. Los modelos de lenguaje clasifican
bien y calibran mal los números: si se les solicita directamente ``asigne una
calificación del 1
al 5'', la misma conversación obtiene calificaciones distintas en dos ejecuciones.
\end{nota}

\subsection{La unidad calificada: la interacción}

Desde el 27 de agosto de 2026 la nota es \textbf{una por interacción}, no una por
conversación. Una conversación de WhatsApp puede contener muchas atenciones de
operadores distintos --- se midieron hasta 14 personas en un mismo hilo --- y una
calificación única las promediaba todas.

Una interacción se corta con las señales del propio CRM:

\begin{center}
\begin{tabular}{@{}lL@{}}
\toprule
\textbf{Señal} & \textbf{Efecto} \\
\midrule
\cod{*resuelto*} & Cierra la interacción. Es la frontera principal \\
\cod{*reabierto*} & Evidencia un cierre que no se aplicó (no corta) \\
Silencio $>$ 6 h & Corta cuando nadie cerró (\cod{SILENCIO\_MAX}) \\
\bottomrule
\end{tabular}
\end{center}

\medskip
Dos reglas de ajuste evitan cortes incorrectos:

\begin{itemize}
\item \textbf{Gracia de cierre: 120 s} (\cod{GRACIA\_CIERRE\_SEG}). El operador cierra
y adjunta el comprobante enseguida; constituye una sola acción y no dos interacciones.
\item \textbf{Cola de cortesía: 600 s} (\cod{GRACIA\_CORTESIA\_SEG}). El ``gracias''
que el cliente envía después del cierre se adjunta a la atención anterior, que fue la
que lo mereció. Sin esta regla, el CRM asigna la conversación a otra persona y esa
persona recibe
\est{1} sin haber escrito una palabra.
\end{itemize}

\begin{nota}
\textbf{El identificador de una interacción es determinista}: se deriva de
(sesión, instante de inicio). Adjuntar la cola de cortesía \emph{absorbe} fragmentos, pero
nunca desplaza un inicio, de modo que recalificar no reasigna identificadores.
\end{nota}

\newpage
% =====================================================================
\section{El orden de precedencia}

Este es el punto que con más frecuencia se malinterpreta: \textbf{no todas las
conversaciones pasan
por el mismo camino}. Existe una cascada de reglas, y prevalece la primera que aplica.

\subsection{Paso 0 --- ¿es evaluable?}

Antes de cualquier rúbrica, un filtro determina si la conversación es calificable.
Las que no, quedan registradas con un motivo de omisión: no desaparecen del tablero,
para que la cobertura sea auditable.

\begin{center}
\begin{tabularx}{\textwidth}{@{}lX@{}}
\toprule
\textbf{Motivo de omisión} & \textbf{Cuándo} \\
\midrule
\cod{grupo\_de\_whatsapp} & Es un grupo. Va \emph{primero}: las rúbricas preguntan si
se le respondió a \emph{un} cliente, y en un grupo no hay uno \\
\cod{internal\_notes\_only} & No hay ningún mensaje real, solo notas del CRM \\
\cod{no\_customer\_reply} & El cliente nunca escribió \\
\cod{customer\_media\_only} & El cliente solo envió imágenes o audio \emph{y} el
operador no resolvió \\
\cod{anomalous\_size} & Más de 250 mensajes reales (\cod{ANOMALOUS\_MESSAGE\_MAX}) \\
\bottomrule
\end{tabularx}
\end{center}

\begin{alerta}
\textbf{Lo que ya NO se omite:} que el negocio no haya contestado. Hasta el 21 de
agosto de 2026 eran 1.167 sesiones omitidas, y esa omisión ocultaba la peor falla que
este sistema puede medir. Medido sobre 300 de ellas, el \textbf{99,7\% tiene una nota
del CRM ``\emph{<Nombre> resuelto la conversación}''}: un operador la cerró sin
escribirle nunca al cliente. Se trata de una acción deliberada y atribuible, no de un
descuido.
\end{alerta}

\subsection{Paso 1 --- la cascada de rúbricas}

\begin{center}
\begin{tabularx}{\textwidth}{@{}clX@{}}
\toprule
\textbf{\#} & \textbf{Regla} & \textbf{Qué decide} \\
\midrule
1 & Sin respuesta      & Nadie del negocio le escribió al cliente $\rightarrow$ \est{1} \\
2 & Cola de cortesía   & El fragmento es solo un ``gracias'' $\rightarrow$ se omite \\
3 & Solo cortesía      & El cliente no planteó nada $\rightarrow$ \est{3} o \est{4} \\
4 & Segmento agente    & La cola es de agentes $\rightarrow$ rúbrica propia \\
5 & Redirección        & La respuesta fue solo un traspaso \\
6 & Por motivo         & El camino general (determinista o LLM) \\
\bottomrule
\end{tabularx}
\end{center}

\begin{nota}
\textbf{Por qué ``sin respuesta'' va primero.} Es el único hecho que no requiere
interpretación: o hay un mensaje del negocio al cliente, o no lo hay. Existe además una
salvedad: \textbf{la nota interna del CRM no cuenta como respuesta}. Es \cod{from\_me}
pero no
es un mensaje al cliente, y contarla convertiría precisamente estas sesiones en ``sí
respondió''.
\end{nota}

\newpage
% =====================================================================
\section{Las rúbricas deterministas}

Varios motivos se califican \textbf{sin LLM}: el resultado se deriva de hechos medibles
(hubo comprobante, hubo confirmación, cuánto tardó). Son enteramente reproducibles y no
consumen inferencia.

\subsection{Depósito}
\emph{Módulo:} \cod{src/deposito.py}\quad\emph{Umbrales:} \cod{AGIL} = 1 min,
\cod{ACEPTABLE} = 5 min

\begin{center}
\begin{tabularx}{\textwidth}{@{}cX@{}}
\toprule
\textbf{Nota} & \textbf{Condición} \\
\midrule
\est{1} & Mandó el comprobante y \textbf{nadie le respondió} \\
\est{2} & Respondió pero \textbf{nunca confirmó} que el dinero se acreditó \\
\est{2} & Confirmó, pero el primer aviso tardó más de 5 min \\
\est{3} & Confirmó, pero el primer aviso tardó más de 1 min \\
\est{4} & Avisó $\le$ 1 min y confirmó, pero cerró sin preguntar si faltaba algo \\
\est{4} & Avisó, confirmó y preguntó --- pero el cliente escribió al final y quedó colgado \\
\est{5} & Avisó $\le$ 1 min, confirmó la acreditación \textbf{y} preguntó si faltaba algo \\
\bottomrule
\end{tabularx}
\end{center}

\paragraph{Dos ramas con techo en \est{4}.} Existen porque el manual de ATC prohíbe
determinados pasos, y penalizarlos equivaldría a exigir el paso prohibido:

\begin{center}
\begin{tabularx}{\textwidth}{@{}lcX@{}}
\toprule
\textbf{Rama} & \textbf{Objetivo} & \textbf{Notas} \\
\midrule
Rechazo (no se podía acreditar) & 1 min & \est{4} si avisó a tiempo, \est{3} si tardó \\
Derivación al agente & 5 min & \est{4} si avisó a tiempo, \est{3} si tardó \\
\bottomrule
\end{tabularx}
\end{center}

\begin{nota}
\textbf{Por qué la derivación tiene 5 minutos y no 1.} El manual (cap. 05) exige que,
antes de derivar, el operador pida el usuario y verifique en el sistema a qué agencia
pertenece. Eso es una consulta, no un reflejo. Exigirle el minuto equivaldría a penalizar la
verificación que el propio manual le ordena realizar.

\textbf{Por qué el techo es \est{4} y no \est{5}.} El 5 significa ``el mejor escenario
del motivo'', y una recarga que no se pudo hacer no lo es.
\end{nota}

\subsection{Retiro}
\emph{Módulo:} \cod{src/retiro.py}\quad\emph{Umbrales:} \cod{AGIL} = 1 min,
\cod{RESPUESTA\_TOPE} = 5 min, \cod{ENTREGA\_AGIL} = 15 min, \cod{ENTREGA\_TOPE} = 30 min

El retiro se mide con \textbf{dos tiempos}: la demora en responder y la demora en
entregar el comprobante de la transferencia.

\begin{center}
\begin{tabularx}{\textwidth}{@{}cX@{}}
\toprule
\textbf{Nota} & \textbf{Condición} \\
\midrule
\est{1} & Pidió el retiro y \textbf{nadie le respondió} \\
\est{2} & Respondió pero \textbf{nunca envió el comprobante} \\
\est{2} & Entregó, pero entrega $>$ 30 min o respuesta $>$ 5 min \\
\est{3} & Entregó, pero respuesta $>$ 1 min o entrega $>$ 15 min \\
\est{4} & Respondió $\le$ 1 min y entregó $\le$ 15 min, sin preguntar si faltaba algo \\
\est{4} & Todo bien, pero el cliente escribió al final y quedó colgado \\
\est{5} & Respondió $\le$ 1 min, entregó $\le$ 15 min \textbf{y} preguntó si faltaba algo \\
\bottomrule
\end{tabularx}
\end{center}

Tiene la misma rama de \textbf{derivación al agente} con techo en \est{4} y objetivo
de 5 minutos, por la misma razón del manual.

\subsection{Registro}
\emph{Módulo:} \cod{src/registro.py}\quad\emph{Umbral:} \cod{ENTREGA\_AGIL} = 5 min

El tiempo se mide desde que el cliente entrega sus datos hasta que recibe su usuario y
contraseña.

\begin{center}
\begin{tabularx}{\textwidth}{@{}cX@{}}
\toprule
\textbf{Nota} & \textbf{Condición} \\
\midrule
\est{1} & Entregó sus datos y \textbf{nadie le respondió} \\
\est{2} & Entregó sus datos pero \textbf{nunca recibió su usuario y contraseña}: el alta quedó a medias \\
\est{3} & Entregó las credenciales, pero demoró más de 5 min \\
\est{3} & Entregó las credenciales, pero el tiempo no es medible (se enviaron antes que los datos) \\
\est{4} & Creó la cuenta $\le$ 5 min, sin acompañar hasta la primera recarga \\
\est{5} & Creó la cuenta \textbf{y además} logró que recargara en la misma conversación \\
\bottomrule
\end{tabularx}
\end{center}

\paragraph{Rama de alta imposible} (el cliente ya tenía cuenta): \est{4} si avisó
dentro de 5 min, \est{3} si tardó. Techo en \est{4} por la misma lógica que las otras
ramas de rechazo.

\subsection{Promoción}
\emph{Módulo:} \cod{src/promo.py}\quad\emph{Umbrales:} \cod{AGIL} = 1 min,
\cod{RAZONABLE} = 5 min, \cod{TOLERABLE} = 15 min

\begin{center}
\begin{tabularx}{\textwidth}{@{}cX@{}}
\toprule
\textbf{Nota} & \textbf{Condición} \\
\midrule
\est{1} & Preguntó por la promo y \textbf{nadie le respondió} \\
\est{2} & Respondió después de 15 min: la consulta perdió vigencia \\
\est{3} & Respondió entre 1 y 15 min \\
\est{4} & Respondió $\le$ 1 min \\
\est{5} & Envió el \textbf{material} de la promoción y respondió $\le$ 5 min \\
\bottomrule
\end{tabularx}
\end{center}

\subsection{Información}
\emph{Módulo:} \cod{src/info.py}\quad\emph{Umbrales:} \cod{AGIL} = 1 min,
\cod{TOLERABLE} = 5 min

\begin{center}
\begin{tabularx}{\textwidth}{@{}cX@{}}
\toprule
\textbf{Nota} & \textbf{Condición} \\
\midrule
\est{1} & Preguntó y \textbf{nadie le respondió} \\
\est{2} & Respondió después de 5 min \\
\est{3} & Respondió entre 1 y 5 min \\
\est{4} & Respondió $\le$ 1 min, sin preguntar si faltaba algo \\
\est{4} & Respondió y preguntó, pero el cliente quedó colgado \\
\est{5} & Respondió $\le$ 1 min \textbf{y} preguntó si faltaba algo \\
\bottomrule
\end{tabularx}
\end{center}

\subsection{Soporte de cuenta}
\emph{Módulo:} \cod{src/soporte.py}\quad\emph{Umbrales:} \cod{AGIL} = 1 min,
\cod{TOLERABLE} = 5 min

Soporte es el motivo con más ida y vuelta (4,5 turnos de media), así que \textbf{no
se mide la primera respuesta sino la mediana de todas las esperas}. Una sola demora
no define una atención de ocho turnos.

\begin{center}
\begin{tabularx}{\textwidth}{@{}cX@{}}
\toprule
\textbf{Nota} & \textbf{Condición} \\
\midrule
\est{1} & Reportó el problema y \textbf{nadie le respondió} \\
\est{2} & Contestó, pero el cliente \textbf{no obtuvo nada concreto}: ni un paso a
seguir ni la certeza de que su caso se escaló \\
\est{2} & Hizo algo por el caso, pero la espera mediana superó los 5 min \\
\est{3} & Atendió el caso, con espera mediana entre 1 y 5 min \\
\est{4} & Atendió rápido y dio una salida concreta, sin preguntar si faltaba algo \\
\est{5} & Atendió rápido, dio una salida concreta \textbf{y} preguntó si faltaba algo \\
\bottomrule
\end{tabularx}
\end{center}

\subsection{Redirección}
\emph{Módulo:} \cod{src/redireccion.py}

Cuando la única respuesta del negocio fue un traspaso a otra línea. Lo que se
califica es \textbf{el destino de la derivación}, no el hecho de derivar.

\begin{center}
\begin{tabularx}{\textwidth}{@{}cX@{}}
\toprule
\textbf{Nota} & \textbf{Condición} \\
\midrule
omisión & Traspaso correcto a una línea propia y activa: no se califica \\
\est{4} & Derivó a otra línea nuestra \textbf{activa}, avisándole al cliente \\
\est{2} & Derivó \textbf{sin dejarle una línea}: el número no se resuelve o la línea está caída \\
\bottomrule
\end{tabularx}
\end{center}

\begin{nota}
Este motivo \textbf{nunca se le pregunta al modelo}. Que la respuesta haya sido solo
un traspaso, que el destino sea una línea nuestra y que esa línea esté activa son
tres hechos que no se pueden verificar leyendo el chat: el último vive en la tabla de
conexiones. \emph{Cuando el hecho es verificable por nosotros, no se consulta al modelo.}
\end{nota}

\subsection{Agilidad (segmento agente)}
\emph{Módulo:} \cod{src/agilidad.py}\quad\emph{Umbrales:} \cod{AGIL} = 1 min,
\cod{BUENO} = 5 min, \cod{ACEPTABLE} = 15 min

Para los agentes cuyo motivo no corresponde a depósito, retiro ni información se mide la
\textbf{peor espera} de todos sus pedidos dentro del horario de atención.

\begin{center}
\begin{tabularx}{\textwidth}{@{}cX@{}}
\toprule
\textbf{Nota} & \textbf{Condición (peor espera)} \\
\midrule
\est{1} & Algún pedido quedó \textbf{sin respuesta} \\
\est{2} & Más de 15 min \\
\est{3} & Entre 5 y 15 min \\
\est{4} & Entre 1 y 5 min \\
\est{5} & Hasta 1 min \\
\bottomrule
\end{tabularx}
\end{center}

\subsection{Sin respuesta y solo cortesía}

\begin{center}
\begin{tabularx}{\textwidth}{@{}lcX@{}}
\toprule
\textbf{Rúbrica} & \textbf{Nota} & \textbf{Condición} \\
\midrule
Sin respuesta  & \est{1} & El cliente escribió y \textbf{nadie del negocio} le escribió \\
Solo cortesía  & \est{4} & El cliente no planteó nada y no quedó colgado \\
Solo cortesía  & \est{3} & El cliente no planteó nada pero se quedó con la última palabra \\
\bottomrule
\end{tabularx}
\end{center}

\newpage
% =====================================================================
\section{El camino del LLM: piso, uplift y modulador}

Los motivos sin rúbrica determinista --- y los casos donde la determinista no
aplica --- pasan por el modelo. Pero el modelo \textbf{no pone la nota}: responde
preguntas sobre hechos y el código deriva la etiqueta.

\subsection{Las tres dimensiones}

\begin{center}
\begin{tabularx}{\textwidth}{@{}llX@{}}
\toprule
\textbf{Dimensión} & \textbf{Rol} & \textbf{Qué mide} \\
\midrule
\cod{resolucion} & piso   & ¿Atendió el motivo de \emph{esta} consulta? \\
\cod{iniciativa} & uplift & ¿Realizó alguna acción más allá del trámite? \\
\cod{cortesia}   & uplift & ¿Saludó, personalizó y utilizó el nombre del cliente? \\
\bottomrule
\end{tabularx}
\end{center}

\subsection{Las reglas, en orden}

\begin{center}
\begin{tabularx}{\textwidth}{@{}Xc@{}}
\toprule
\textbf{Situación} & \textbf{Etiqueta} \\
\midrule
Maltrato grave & mala (\est{1}) \\
No atendió \textbf{y} hubo fricción (el cliente insistiendo, sin respuesta) & mala (\est{1}) \\
No atendió & deficiente (\est{2}) \\
Atendió, pero hubo \textbf{fricción} (reinsistió sin respuesta) & deficiente (\est{2}) \\
Atendió, pero fue \textbf{confuso} \emph{y} está corroborado & deficiente (\est{2}) \\
Atendió limpio \textbf{+} acción extra \emph{o} cortesía destacada & excelente (\est{5}) \\
Atendió limpio & buena (\est{4}) \\
\bottomrule
\end{tabularx}
\end{center}

\begin{nota}
\textbf{Por qué hacer bien el trabajo ya vale \est{4}.} Hasta la versión 3, el piso
limpio se limitaba a \est{3} y para superarlo se requería empuje comercial. Medido
sobre 213 sesiones de depósito: 149 respondieron en $\le$ 2 minutos \emph{y}
confirmaron la acreditación --- el trabajo completo --- y 135 de esas quedaron en
\est{3}. Ejecutarlo de forma impecable representaba $+0{,}13$ estrellas frente a no
hacerlo. La
escala no medía lo que declaraba medir.
\end{nota}

\subsection{Por qué el código corrige al modelo}

El sistema aplica \textbf{overrides deterministas} sobre lo que responde el LLM,
en los dos sentidos:

\begin{itemize}
\item Si el modelo dice ``no atendió'' pero el operador \textbf{resolvió} o
\textbf{empujó} de forma verificable, se corrige a ``atendió''.
\item Si el modelo dice ``no atendió'' pero \textbf{el propio cliente confirmó} que
se resolvió, se corrige. El cliente es el único testigo que sabe si su problema se
resolvió.
\item Si el modelo dice ``maltrato grave'' pero la señal determinista no lo
encuentra, se corrige.
\item Un ``confuso'' del modelo \textbf{sin corroboración} (el cliente ni preguntó ni
reinsistió) no reduce la calificación: se limita a \est{3}.
\end{itemize}

\newpage
% =====================================================================
\section{El reloj: qué cuenta como espera}

\subsection{Horario de atención}

Todos los tiempos se miden \textbf{descontando la noche}. El horario es de
\textbf{06:00 a 23:59} hora de Ecuador (18 horas por día).

\begin{alerta}
\textbf{Esto no es un detalle.} Un cliente que escribe a las 00:06 y recibe respuesta
a las 06:05 no esperó seis horas: esperó \textbf{cinco minutos de horario}. Publicar
el tiempo de reloj corrido imputaría negligencia a quien respondió a los seis minutos
de iniciado el turno. La misma función que descuenta la noche la usan todas las rúbricas.
\end{alerta}

\subsection{Precisión}

Los tiempos se comparan \textbf{redondeados al segundo}. Los objetivos del manual están
expresados en minutos enteros y se comparaban contra marcas de tiempo de WhatsApp con
milisegundos: dos operadores perdieron dos estrellas por 365 y 312 \emph{milisegundos}.

\newpage
% =====================================================================
\section{Omisiones: lo que no se califica y por qué}

\begin{center}
\begin{tabularx}{\textwidth}{@{}lX@{}}
\toprule
\textbf{Motivo} & \textbf{Razón} \\
\midrule
\cod{grupo\_de\_whatsapp} & No hay \emph{un} cliente a quien se le responda \\
\cod{internal\_notes\_only} & Solo notas del CRM, ningún mensaje real \\
\cod{no\_customer\_reply} & El cliente nunca escribió \\
\cod{customer\_media\_only} & Solo media, sin texto, y el operador no resolvió \\
\cod{anomalous\_size} & Más de 250 mensajes: patológico \\
\cod{redireccion} & Traspaso limpio a una línea propia y viva \\
\cod{datos\_incompletos} & El alta se abandonó porque el cliente no pasó sus datos \\
\cod{cola\_de\_cortesia} & El fragmento es solo el ``gracias'' de una atención ya cerrada \\
\bottomrule
\end{tabularx}
\end{center}

\begin{nota}
\textbf{La omisión por cola de cortesía (v24).} Decisión del negocio, textual:
\emph{``merece el skip porque para el usuario es la misma conversación, no quiere nada
más; no se debe castigar a ATC por algo que no requiere atención''}.

Funciona en dos capas: una determinista y gratuita, y para el residuo que la primera
no reconoce, el modelo --- con un costo medido de \textbf{0,8 inferencias por día}.
Ante cualquier falla del modelo se puntúa: nunca se pierde un reclamo por una
inferencia que no se completó.
\end{nota}

\newpage
% =====================================================================
\section{Limitaciones conocidas}

Este sistema tiene deudas medidas y documentadas. Omitirlas restaría utilidad a este
documento.

\begin{enumerate}
\item \textbf{El detector de acreditación está calibrado a la plantilla vieja.}
Un operador que cierra con ``¡Está listo!'' o ``Lista su recarga'' no es reconocido.
Medido sobre 2.784 depósitos: quien \emph{no} usa esas frases promedia 4,00 y tiene
4,1\% de notas $\le$ 2; quien las usa promedia 3,04 y tiene \textbf{48,1\%}. Es
\textbf{11,7 veces} más probable sacar una nota baja por usar la frase propia. No se
trata de
la conducta del operador, sino del detector.

\item \textbf{La quinta estrella es inalcanzable en Facebook.} El requisito de
preguntar ``¿algo más?'' no discrimina por motivo ni por antigüedad: discrimina por
\textbf{canal}. Sobre 162 atenciones en Facebook, \textbf{0} lo cumplen, porque las
respuestas rápidas que contienen esa pregunta no existen en ese canal.

\item \textbf{El segmento no entra en el scoring.} Un agente --- un vendedor
profesional --- se califica con la misma vara comercial que un jugador. Existe una
rúbrica propia para agentes, pero solo para depósito, retiro, información y agilidad.

\item \textbf{Las sesiones que el CRM reabre no se parten.} Se midió: partirlas
dejaría al 18\% sin evaluar y perjudicaría a 25 casos concretos. Se trata de una decisión
deliberada, no de una omisión.
\end{enumerate}

% =====================================================================
\section{Dónde vive cada cosa}

\begin{center}
\begin{tabularx}{\textwidth}{@{}lX@{}}
\toprule
\textbf{Módulo} & \textbf{Responsabilidad} \\
\midrule
\cod{src/router.py}        & Decide si la conversación es evaluable \\
\cod{src/interacciones.py} & Corta la conversación en interacciones \\
\cod{src/rubrics.py}       & Escala, dimensiones y derivación de la etiqueta \\
\cod{src/scorer.py}        & Camino del LLM y overrides deterministas \\
\cod{src/horario.py}       & Horario de atención y espera efectiva \\
\cod{src/signals.py}       & Señales deterministas sobre el transcript \\
\cod{src/worker.py}        & Cascada de precedencia y persistencia \\
\cod{src/deposito.py} \dots & Una rúbrica determinista por motivo \\
\bottomrule
\end{tabularx}
\end{center}

\vfill
\begin{center}
\small\color{gray70}
Versión del scoring: \cod{2026.08-rubricas-v24} ---
el registro completo de cambios vive en \cod{src/store.py}.
\end{center}

\end{document}
